\setchapterpreamble[u]{\margintoc}
\chapter{Measurement of the Transfer Function}
\labch{tf-measurement}

If we have an optical system and want to model the way light passes through that system, either in order to enhance our
    understanding of the structure (source) that causes the signal or for any other reason, we need to determine one of
    the transfer functions introduced in the previous part. % TODO: reference <14-11-23> 
Of the transfer functions, it is most intuitive to directly measure the \ac{PSF} or more specifically the
    magnitude of the \ac{PSF}.
We disregard the phase change that the optical system may introduce, because in most set-ups it is actually impossible 
    to directly measure the imaginary part of the \ac{PSF} since the light sensors\sidenote{In our case the \ac{CMOS} 
    sensor.} can only measure the energy that hits the sensor which is the same for all phases of the \ac{EM} wave 
    since it only depends on the intensity of the \ac{EM} wave as you may recall from \refch{light}. 
    % TODO: Add closer reference. <14-11-23> 
\marginnote{It is possible in some cases to determine algorithmically the imaginary part (or the phase) from the magnitude
    algorithmically which will be discussed in \refsec{phase-retrieval}.}

\section{Acquisition Model}%
\labsec{acquisition-model}

% TODO: First step, noise. Measure noise from the background of the beads image. <14-11-23> 
First off, let us devise an acquisition model based on the way, we expect the light from the source to manifest itself
    in the resulting signal that is captured by the sensor.


\subsection{Illumination}%
\labsubsec{illumination}
To determine the transfer function of the optical system, means to measure it for a uniformly illuminated focal plane.
However, we do not have data directly from a simple uniformly illuminated wide field system.
Our images are acquired with the harmonic patterns discussed in the previous section. % TODO: Add reference <19-11-23> 
This means that for a still\sidenote{Meaning that the source of the light is stationary in contrast to \emph{in vivo}
    biological acquisitions, where this is not the case.} image, we have 9 wide field images with 3 different pattern 
    orientations and 3 different phases\sidenote{In accordance with the method that is used to generate the patterns in 
    the case of our \ac{SIM} setup, we are assuming that the phases differ by \(2\pi/3\) from one another and do not 
    determine them independently.} for each of those.

\begin{figure}[h]
% \begin{marginfigure}
    \png{beads-9-patterns-by-orientations-and-phases}
    % \marginpng{beads-9-patterns-by-orientations-and-phases}
    \caption{Acquisitions of 9 wide field images with harmonic patterned illuminations. The orientations are labelled
    1, 2 and 3 and the phases \(-\), \(0\) and \(+\).}
    \labfig{beads-LR-9-images}
% \end{marginfigure}
\end{figure}

In order to account for this discrepancy, we will imitate the uniform illumination by summing up all the images acquired
    with different illuminations together.
\marginnote{This will also reduce noise, which will become evident in \refsubsec{noise}.}
Due to additivity of the optical system, this results in the same signal sum, as if an acquisition was made with a sum
    of the illumination patterns.
In one orientation, without loss of generality, we assume that the angle of the illumination harmonic vector with the
    \(x\)-axis is \(0\) and for \(y = 0\) we have

\begin{marginfigure}
  \marginpdftex{sum-of-patterns-z-x-axes-for-one-orientation}
  \caption{Sum of the illumination pattern intensities results in a uniform intensity.}
  \labfig{sum-of-ips}
\end{marginfigure}

\begin{equation*}
    \begin{split}
        I_\Sigma ( \pmqty{0 \\ y}) &= \sum_{i=1}^3 1 - \frac{m_i}{2} \cos\ab(2\pi \pmqty{0 & k_y} \cdot \pmqty{0 \\ y}
        + \frac{2\pi i}{3}) \\
                                    & = 3 - \sum_{i=1}^3 \frac{m}{2} \cos\ab(2\pi k_y y + 
                                    \frac{2\pi i}{3}) \\
                                    & = 3 -\frac{m}{2} \cos\ab(2\pi k_y y)\underbrace{\sum_{i=1}^3\ab(\cos\frac{2\pi
                                    i}{3})}_{=0}
                                    - \sin\ab(2\pi  k_y y)\underbrace{\sum_{i=1}^3\ab(\sin\frac{2\pi i}{3})}_{=0} \\
    \end{split}
\end{equation*}

% TODO: Section about the images <16-11-23> 
% TODO: Figure: Add reconstructed (HR) image figure and LR image figures <16-11-23> 

% TODO: Second step, illumination... The sum <16-11-23> 
    % TODO: Side figure of the z-x axis cut of the illumination intensity cut <16-11-23> 
    % TODO: Note about additivity and why the 3× illuminated signal is the same as if we would have illuminated it all at
    % once <16-11-23> 

\subsection{Noise}%
\labsubsec{noise}

As a first step towards developing a faithful model, we need to include noise \(\eta(\vec{x})\) in the model.

The mode of noise that we will be particularly interested in is \emph{additive noise}, in which the strength of the 
    noise is added to the noise free signal:
\begin{equation}
    \labeq{additive-noise-model}
    u(\vec{x}) = f(\vec{x}) + \eta(\vec{x}).
\end{equation}
In \refeq{additive-noise-model} the noise \(\eta(\vec{x})\) can have a specified distribution\sidenote{It is not 
    uncommon to choose to model the noise by a Gaussian distribution (\(\distribution{D} \equiv \distribution{N}\)) even
    if the Gaussian distribution is physically unrealistic for the given noise source. This simplification is often made 
    for the purpose of simplifying computation.}
\begin{equation*}
    \eta(\vec{x}) \sim \distribution{D}.
\end{equation*}
An important property to consider when characterising noise from different sources is \emph{signal independence}, i.e. 
    whether the strength of the noise is independent from the intensity of the signal
\begin{equation*}
    \left\{ \eta(\vec{x}) \mid f(\vec{x})  \right\} = \eta(\vec{x}).
\end{equation*}

There are many sources of noise that should be considered in the acquisition model.
The most prominent of those sources are caused by the image sensor (\ac{CMOS}). %  and the fluorescent proteins (\ac{EGFP}).
It is also common to regard the variations in the ambient (background) illumination i.e. the light that is not part of 
    the source we are measuring as noise.
The most prominent sources of noise, that are burdening the expected signal are:
\begin{description}
    \item[Photon noise]
        Inherent to the quantized nature of energy transmission of light. 
        Sensors measure irradiance by counting the number of discrete photons that land on the sensor in a given time
            interval. 
        The temporal distribution of the independent arrivals of individual photons can treated as a Poisson process.
        The number \(N\) of photons detected in the time interval \(\Delta\tau\) can thus be described by the
            distribution
            \begin{equation*}
                \distribution{P}\left[ N = k \right] = \frac{e^{-\lambda\cdot\Delta\tau}(\lambda\cdot\Delta\tau)^k}{k!},
            \end{equation*}
            where the rate parameter \(\lambda\cdot\Delta\tau\) corresponds to the expected incident photon count, i.e.
            the expected number of photon detections based on the source luminance.
        Since \(\E N = \Var(N) = \lambda\cdot\Delta\tau\), the ratio of signal to noise increases with
            \(\sqrt{\lambda\cdot\Delta\tau}\) and so this source of noise is most prominent in
            situations with low signal levels.
        For larger counts of photons, we can use the central limit theorem to approximate the photon noise by a Gaussian
            additive noise \sidecite{hasinoff2021}, whose variance and mean both depend on the signal level
            \begin{gather*}
                u(\vec{x}) = f(\vec{x}) + \eta(\vec{x}) \\
                \eta(\vec{x}) \sim \distribution{N}\left(\lambda(\vec{x})\cdot\Delta\tau, 
                                                         \lambda(\vec{x})\cdot\Delta\tau \right).
            \end{gather*}
    \item[Dark current]
        Semiconductor detectors in the image sensor are susceptible to thermal agitation, which can cause the
            photocapacitors to charge even in the absence of light (hence the name dark current) leading to false
            signals.
            Due to the phenomenon of dark current, the sensors are often cooled to reduce this agitation.
            Noise from this source is signal independent and additive.
    \item[Background noise]
        Additive and independent noise from the ambient light that reaches the sensor or unwanted light caused by other
        effects\sidenote{An example can be the lagged emission of the fluorescent proteins.}.
        It is most prominent for low energy light, that is light with low frequency (e.g. infrared light).
    \item[Auxiliary sources]
        A variety of causes mainly in the imaging sensor that are produced at the read-out time\sidenote{The phases
            after the photo activation of the sensor elements (consisting of amplification,
            digitization of the analog signal etc.).}.
        Notably the charge transfer in the \ac{CMOS} sensor can be the cause noise.
        Auxiliary sources can be either signal dependent or independent.
\end{description}

If we know or assume a distribution of an additive noise, it is possible to estimate its parameters using parts of the
    image with approximately or theoretically constant intensity.
Importantly, it is often the case that we assume or in approximation assume that the distribution of the additive noise 
    is Gaussian and we can estimate the mean and the standard deviation of the noise using the \ac{MLE}
\begin{gather*}
    A \coloneq \left\{ \vec{x} \mid f(\vec{x}) \approx c  \right\} \\
    \hat{\mu}_{\eta} = \frac{1}{|A|} \sum_{\vec{x} \in A} (u(\vec{x}) - c) \\
    \hat{\sigma^2}_{\eta} = \frac{1}{|A|} \sum_{\vec{x} \in  A} (u(\vec{x}) - \hat{\mu}_{\eta})^2.
\end{gather*}
Furthermore even if we do not want to limit the theorized distribution of noise, we can eyeball or even make an
    empirical estimate of the distribution using the histogram of the intensities in \(A\).
In our case, the premise that allows us to select such a set \(A\) is that besides the beads, there are no other sources
    of signal.
The problem therefore reduces to the distinction of the beads from the background for which an Otsu
    threshold\sidecite{otsu1979} can be used.

% TODO: Figure: Add background histogram figure <16-11-23> 


% TODO: Third step, transfer function. <14-11-23> 
% TODO: Fourth step, signal. <14-11-23> 

% \section{Setting}%
% \labsec{setting}
%
% \section{PSF Model}%
% \labsec{psf-model}
%
% \section{Locating the Beads}%
% \labsec{locating-the-beads}
%
% \section{Estimating the PSF from the Model}%
% \labsec{estimating_the_psf_from_the_model}

\section{Phase Retrieval}%
\labsec{phase-retrieval}

%
% \subsection{Error Estimation}%
% \labsubsec{error_estimation}
